\subsection{Modelagem dinâmica}
\label{sec:MR_modelagem_dinamica}

O comportamento dinâmico é obtido via Lagrangiana, pois evita a enumeração explícita de todas as forças de ligação.
A Lagrangiana é dada por:
\begin{equation}
\mathcal{L}(\vec q,\dot{\vec q}\,) \;=\; T(\vec q,\dot{\vec q}\,) \;-\; V(\vec q\,),
\end{equation}

sendo a energia cinética expressa por:
\begin{equation}
\label{eq:energia_cinetica_lagrangiana}
T(\vec q,\dot{\vec q}\,)
\;=\;
\tfrac12\,\dot{\vec q}\,^T\,\mathbf M(\vec q\,)\,\dot{\vec q}.
\end{equation}

Neste trabalho, como o manipulador possui quatro juntas revolutas e uma prismática, tem-se 
$\vec q \in \mathbb{R}^{5}$, de modo que a matriz de inércia generalizada 
$\mathbf M(\vec q)$ é simétrica, definida positiva e de dimensão $5\times 5$.

\subsubsection{Coordenadas generalizadas}
As coordenadas generalizadas são expressas por:

\begin{equation}
\vec q
\;=\;
\begin{bmatrix}\theta_1 & \theta_2 & \theta_3 & \theta_4 & d_5\end{bmatrix}^T,
\end{equation}

sendo velocidades $\dot{\vec q}$ e
as acelerações $\ddot{\vec q}$.

\subsubsection{Equações de Euler--Lagrange}
Para cada coordenada generalizada $q_i$, tem-se:

\begin{equation}
\frac{d}{dt}\!\left(\frac{\partial \mathcal L}{\partial \dot q_i}\right)
\;-\; \frac{\partial \mathcal L}{\partial q_i}
\;=\;
Q_i, \qquad i=1,\dots,n,
\label{eq:EL_componentes}
\end{equation}
em que $Q_i$ é a força generalizada associada a $q_i$.

\subsubsection{Forma matricial compacta}
A partir de \eqref{eq:energia_cinetica_lagrangiana} e de $V(\vec q\, )$, obtém-se:

\begin{equation}
\label{eq:EL_compacta}
\mathbf M(\vec q\,)\,\ddot{\vec q} \;+\; \mathbf C(\vec q,\dot{\vec q}\,)\,\dot{\vec q} \;+\; \vec g(\vec q\,)
\;=\;
\vec \tau,
\end{equation}

onde $\vec\tau$ são os torques ou forças exercidos pelos atuadores,
$\vec g(\vec q)$ é o termo gravitacional e
$\mathbf C(\vec q,\dot{\vec q}\,)\,\dot{\vec q}$ reúne os efeitos de Coriolis e centrífugos.

\subsubsection{Termo gravitacional e microgravidade}
As coordenadas generalizadas em função da gravidade, pode ser expressa por:

\begin{equation}
\vec g(\vec q\,)
\;=\;
\left(\frac{\partial V(\vec q\,)}{\partial \vec q}\right)^T.
\label{eq:g_def}
\end{equation}

No presente trabalho, o potencial gravitacional foi desconsiderado, pois as operações são em ambiente de microgravidade, reduzindo \eqref{eq:EL_compacta} a:

\begin{equation}
\label{eq:EL_microgravidade}
\mathbf M(\vec q\,)\,\ddot{\vec q}
\;+\;
\mathbf C(\vec q,\dot{\vec q}\,)\,\dot{\vec q}
\;=\;
\vec \tau.
\end{equation}

\subsubsection{Forças generalizadas força--momento}
Tão logo a força--momento atua na ferramenta, a contribuição nas juntas é dada por:

\begin{equation}
\label{eq:tau_Jt_w}
\vec \tau
\;=\;
\mathbf J^{T}(\vec q)\,\vec w_{E}, 
\qquad
\vec w_{E} \;=\; \begin{bmatrix}\vec f_{E} \\ \vec m_{E}\end{bmatrix},
\end{equation}

em que $\mathbf J(\vec q)$ é a Jacobiana geométrica
que mapeia $\dot{\vec q}$ para a torção
$[\vec v_{E}; \vec \omega_{E}]^T$.

\subsubsection{Energia mecânica e balanço de potência}
Para realizar uma verificação numérica e análise de estabilidade, pode-se multiplicar
\eqref{eq:EL_compacta} à esquerda por
$\dot{\vec q}^T$, e obtém-se:

\begin{equation}
\frac{d}{dt}\!\left(
    \tfrac12\,\dot{\vec q}^{\,T}\mathbf M(\vec q\, )\dot{\vec q} + V(\vec q\, )\right)
\;=\;
\dot{\vec q}\, ^T\,\vec \tau.
\label{eq:energy_balance}
\end{equation}

%==================================================
% Energia cinética por elo e matriz de inércia
%==================================================
\subsubsection{Energia cinética do MR}
\label{sec:MR_energia_cinetica}

Cada elo $i$ possui massa $m_i$, centro de massa $C_i$ e tensor de inércia ${}^{0}\!\mathbf I_{C_i}$, e sua energia cinética é dada por:

\begin{equation}
\label{eq:energia_elos}
T_i
\;=\;
\tfrac12\, m_i \,\vec v_{C_i}^{\,T}\vec v_{C_i}
\;+\; \tfrac12\, \vec \omega_i^{\,T}\, ({}^{0}\!\mathbf I_{C_i}) \, \vec \omega_i,
\end{equation}

sendo que a velocidade no centro de massa é:
\begin{equation}
\vec v_{C_i}
=
\mathbf J_{v,C_i}(\vec q\,)\,\dot{\vec q}
\end{equation}

enquanto que a velocidade angular é expresso por:
\begin{equation}
\vec \omega_i
=
\mathbf J_{\omega,i}(\vec q\, )\,\dot{\vec q}
\end{equation}

Portanto, a energia cinética total será:
\begin{equation}
\label{eq:T_soma_links}
T(\vec q,\dot{\vec q})
\;=\;
\sum_i \left(\tfrac12\, m_i\,\|\mathbf J_{v,C_i}(\vec q)\,\dot{\vec q}\|^2 \;+\; \tfrac12\, \dot{\vec q}^{\,T}\mathbf J_{\omega,i}^{T}(\vec q)\,{}^{0}\!\mathbf I_{C_i}\,\mathbf J_{\omega,i}(\vec q)\,\dot{\vec q}\right).
\end{equation}

Resultando assim, na matriz de inércia generalizada:

\begin{equation}
\label{eq:definicao_Mq}
\mathbf M(\vec q)
\;=\;
\sum_i \Big(m_i\,\mathbf J_{v,C_i}^{T}(\vec q\,)\,\mathbf J_{v,C_i}(\vec q\, ) \;+\; \mathbf J_{\omega,i}^{T}(\vec q\, )\,{}^{0}\!\mathbf I_{C_i}\,\mathbf J_{\omega,i}(\vec q\, )\Big),
\end{equation}

que é simétrica e definida positiva:
\begin{equation}
\mathbf M(\vec q\, )=\mathbf M^{T}(\vec q\, ), 
\qquad
\vec x^{T}\mathbf M(\vec q\, )\vec x>0 \ \ (\vec x\neq \vec 0\, ).
\end{equation}